\chapterhead{2}{Round on Purpose}{Choosing a number that says no more than you know}\label{ch2}

Mia's pencil ends exactly halfway between the 17 and 18 marks on her ruler. The ruler has no mark there. She has to choose.

Choosing the nearest mark is called \textbf{rounding}. It is not cheating, and it is not sloppy. It is a way of saying: \emph{this is as close as my ruler lets me get.}

\subhead{Which mark is nearest?}
Every reading sits somewhere between two marks. Rounding picks the nearer one. Look at a number line for lengths between 61 and 62~cm.

\begin{center}
\begin{tikzpicture}[x=9cm]
  \draw[line width=1.2pt] (0,0) -- (1,0);
  \foreach \x/\l in {0/61,0.5/61.5,1/62}{
    \draw[line width=1.2pt] (\x,-0.15) -- (\x,0.15);
    \node[font=\sffamily\bfseries\small,below=4pt] at (\x,-0.15) {\l};
  }
  \foreach \x in {0.1,0.2,0.3,0.4,0.6,0.7,0.8,0.9}{\draw (\x,-0.07) -- (\x,0.07);}
  \fill[sky] (0.2,0) circle (3pt) node[above=6pt,font=\sffamily\small\bfseries] {61.2};
  \fill[tangerine] (0.5,0) circle (3pt) node[above=6pt,font=\sffamily\small\bfseries] {61.5};
  \fill[leaf] (0.8,0) circle (3pt) node[above=6pt,font=\sffamily\small\bfseries] {61.8};
  \draw[-{Stealth},line width=1.1pt,color=sky] (0.2,0.9) -- (0.02,0.9) node[midway,above,font=\sffamily\scriptsize] {rounds to 61};
  \draw[-{Stealth},line width=1.1pt,color=leaf] (0.8,0.9) -- (0.98,0.9) node[midway,above,font=\sffamily\scriptsize] {rounds to 62};
\end{tikzpicture}
\end{center}

61.2 is closer to 61, so it rounds \emph{down} to 61. 61.8 is closer to 62, so it rounds \emph{up}. Exactly halfway, at 61.5, we have to agree on a rule. In this book, \textbf{halves round up}, so 61.5 becomes 62.

\begin{tcolorbox}[colback=sun!14,colframe=sun!85!black,boxrule=1.2pt,arc=2mm,left=6pt,right=6pt,top=4pt,bottom=4pt,before skip=8pt,after skip=8pt]
\sans\bfseries\color{ink} The rounding rule\par
\normalfont\sffamily\small Look at the digit just after the place you want to keep.\\
\textbf{0, 1, 2, 3, 4:} the number stays. \qquad \textbf{5, 6, 7, 8, 9:} the number goes up by one.\\
Example: 38.4 to the nearest whole number is 38. 38.5 is 39. 38.6 is 39.
\end{tcolorbox}

\subhead{Rounding is a pop}
When we round 17.4~cm to 17~cm, we squash a number into a simpler one. That is a pop, so we must say what we kept and what we lost.

\begin{center}
\begin{tikzpicture}
  \node[bigcir,minimum size=2.9cm] (a) at (0,0) {};
  \node[cir] at (0,0.35) {17.4 cm};
  \node[note,text width=2.2cm] at (0,-0.6) {pencil, read with an mm ruler};
  \node[valcir] (b) at (5.2,0) {17 cm};
  \draw[popA] (a) -- node[above,note,text width=2.2cm] {round to the nearest cm} (b);
  \node[note,align=left,anchor=west,text width=3.4cm] at (6.4,0) {\textbf{kept:} about how long\\\textbf{lost:} the 0.4 cm};
\end{tikzpicture}
\end{center}

\trust{We compared 17.4 with the rounded 17.}
{A number rounded to the nearest whole is never more than half a unit away from the original. The gap here is 0.4~cm, which is less than 0.5.}
{Whether the 17.4 reading itself was right, and whether adding up many rounded numbers stays close to the true total. It may not. You will see why in a moment.}

\subhead{Rounding to the nearest ten}
You can round to bigger steps too. To round to the nearest ten, look at the ones digit: 34 becomes 30, and 35 becomes 40. People round a distance of 487~m to ``about 490~m'' or ``about 500~m'', depending on how much detail the listener needs.

\subhead{Rounding in the news}
Suppose a town has exactly 12\,482 people on counting day. Here is what different people might write. Each one is honest, as long as it says ``about''.

\begin{center}
\sffamily\small\renewcommand{\arraystretch}{1.5}
\begin{tabular}{@{}lr@{\qquad}l@{}}
\toprule
\textbf{Rounded to the nearest} & \textbf{Number} & \textbf{Who might use it}\\
\midrule
ten & 12\,480 & the town hall\\
hundred & 12\,500 & the local paper\\
thousand & 12\,000 & a national newspaper\\
ten thousand & 10\,000 & a map of a whole country\\
\bottomrule
\end{tabular}
\end{center}

Every step up throws away more detail, but each one keeps the answer to ``about how big?'' for a different audience.

\begin{multix}[hmm]\multixbot[hmm]\tcblower
The newspaper said 12\,000 people live in the town. So the real number is exactly 12\,000. I am an expert reader of newspapers. Or am I?
\end{multix}

\begin{tryit}{Round and add}
\textbf{You need:} a ruler marked in mm, and five small things such as a pencil, an eraser, a paper clip, a key and a crayon.
\begin{enumerate}[leftmargin=1.6em]
\item Measure each thing to the nearest millimetre, in centimetres (for example 17.4~cm). Write each in the first column.
\item Round each length to the nearest centimetre. Write it in the second column.
\item Add the first column. Then add the second column.
\item Compare the two totals.
\end{enumerate}
\par\medskip
\renewcommand{\arraystretch}{1.7}
\begin{tabularx}{\linewidth}{|>{\sffamily\small}p{2.6cm}|X|X|}
\hline
\textbf{Thing} & \textbf{To the nearest mm (cm)} & \textbf{Rounded to the nearest cm}\\ \hline
 & & \\ \hline
 & & \\ \hline
 & & \\ \hline
 & & \\ \hline
 & & \\ \hline
\textbf{Total} & & \\ \hline
\end{tabularx}
\end{tryit}

\begin{worked}{Practice numbers (made up for practice, not from a real measurement)}
Pencil 17.4, eraser 5.2, paper clip 3.4, key 5.7, crayon 9.1.\\
Exact total: $17.4+5.2+3.4+5.7+9.1=40.8$~cm. \textbf{Rounded to the nearest cm:} 41.\\
Round each first: $17+5+3+6+9=40$. The total of the rounded numbers is \emph{off by 1} from the rounded total of the exact ones.
\end{worked}

Each rounding threw away a little bit, and the little bits added up. That is why careful people round \emph{at the end}, not at every step.

\begin{keepbox}
When we rounded each of the five lengths, \textbf{we kept} \wline[4cm] \textbf{and we set aside} \wline[4cm].
\end{keepbox}

\subhead{How many digits are worth writing?}
Here is a puzzle that fools many grown-ups. You cut a 17~cm piece of string into 3 equal pieces. How long is each piece?

Your calculator says $17\div3=5.6666667$.

But your ruler only has marks every millimetre. The very last digit of 5.6666667 is about a millionth of a millimetre! No ruler can show that, so writing it pretends you know something you cannot know. The honest answer is \textbf{5.7~cm}.

\begin{multix}[hmm]\multixbot[hmm]\tcblower
The calculator gave me 5.6666667, so my string piece is 5.6666667~cm. More digits means more exact. Right? I think that is wrong. Can you say why?
\end{multix}

\begin{center}
\sans\bfseries\large\color{ink}
Do not write more digits than your tool can honestly give.
\end{center}

\clearpage

\begin{puzzle}{Round the number}
Round each number. Remember: halves round up.\medskip

\renewcommand{\arraystretch}{1.8}
\begin{tabular}{@{}l>{\sffamily\bfseries}cl>{\sffamily\bfseries}c@{}}
\multicolumn{2}{@{}l}{\textbf{To the nearest whole number}} & \multicolumn{2}{l}{\textbf{To the nearest ten}}\\
(a) 23.4 & \wline[1.8cm] & (f) 45.2 & \wline[1.8cm]\\
(b) 23.5 & \wline[1.8cm] & (g) 81 & \wline[1.8cm]\\
(c) 7.49 & \wline[1.8cm] & (h) 135 & \wline[1.8cm]\\
(d) 99.6 & \wline[1.8cm] & & \\
(e) 12.51 & \wline[1.8cm] & & \\
\end{tabular}
\end{puzzle}

\begin{puzzle}{Sensible digits}
Your calculator shows a long answer. Your ruler is marked in millimetres. Write each answer in centimetres, to the nearest millimetre.
\begin{center}\sans\small\renewcommand{\arraystretch}{1.8}
\begin{tabular}{@{}lll@{}}
(a) $17\div3=5.6666667$ & write: & \wline[2cm]~cm\\
(b) $33\div7=4.7142857$ & write: & \wline[2cm]~cm\\
(c) $50\div8=6.25$ & write: & \wline[2cm]~cm\\
\end{tabular}
\end{center}
\end{puzzle}

\begin{puzzle}{Which rounding is wrong?}
Four people rounded these lengths to the nearest centimetre. Only one rounding is wrong.
\begin{center}\sans
17.4 $\to$ 17 \qquad 8.5 $\to$ 8 \qquad 26.7 $\to$ 27 \qquad 3.2 $\to$ 3
\end{center}
Which one is wrong? \wline[2.5cm] What should it be? \wline[2cm]
\end{puzzle}

\begin{puzzle}{Ribbons}
Three pieces of ribbon are each 1.4~cm long.\\
(a) What is their total length, exactly? \wline[2cm]~cm\\
(b) Round each piece to the nearest centimetre first, then add. What do you get? \wline[2cm]~cm\\
(c) The two answers are different. Which total is closer to the truth, and why?\\
\rule{\linewidth}{0.5pt}
\end{puzzle}

\begin{puzzle}{Draw the pop}
Fill in the empty circle for each pop. Then write what was kept and what was lost.
\begin{center}
\begin{tikzpicture}[font=\sffamily\small]
  \foreach \y/\from/\rl in {0/8.7 cm/nearest cm,-2.1/46/nearest ten,-4.2/2.35 m/nearest 0.1 m}{
    \node[bigcir,minimum size=1.5cm] (a) at (0,\y) {\bfseries\from};
    \node[valcir,dashed,minimum size=1.5cm] (b) at (5,\y) {};
    \draw[popA] (a) -- node[above,note,text width=2.6cm] {round to the \rl} (b);
    \node[anchor=west,align=left,text width=5.2cm,font=\sffamily\scriptsize] at (6.2,\y) {kept: \wline[3.4cm]\\[5pt]lost: \wline[3.4cm]};
  }
\end{tikzpicture}
\end{center}
\end{puzzle}

\begin{puzzle}{Rounding maze}
Start at the top left. Move only \textbf{right} or \textbf{down}, one square at a time. You may step only on a number that rounds to \textbf{50} when rounded to the nearest ten. Reach the bottom right.
\begin{center}
\begin{tikzpicture}[font=\sffamily\bfseries\large]
  \foreach \r/\c/\n in {0/0/47,0/1/61,0/2/38,1/0/52,1/1/54,1/2/29,2/0/33,2/1/49,2/2/45}{
    \node[draw=ink,line width=1pt,fill=white,minimum size=1.35cm] at (\c*1.35,-\r*1.35) {\n};
  }
  \node[font=\sffamily\bfseries\scriptsize,text=leaf!70!black] at (-1.2,0) {START};
  \node[font=\sffamily\bfseries\scriptsize,text=tomato] at (3.9,-2.7) {END};
\end{tikzpicture}
\end{center}
\end{puzzle}
\clearpage
